\documentclass[a4paper,12pt]{article}
\usepackage[italian]{babel}
\usepackage[utf8]{inputenc}
\usepackage[T1]{fontenc}
\usepackage{geometry}
\usepackage{setspace}
\usepackage{listings}
\usepackage{xcolor}

\geometry{margin=2.5cm}
\singlespacing

\title{Verifica di Intelligenza Artificiale - Classe 4AIQ}
\author{Prof. Fedeli Massimo}
\date{}

\definecolor{codegray}{rgb}{0.95,0.95,0.95}

\lstset{
	backgroundcolor=\color{codegray},
	basicstyle=\ttfamily\small,
	breaklines=true,
	frame=single,
	showstringspaces=false,
	tabsize=4,
	language=Python
}

\begin{document}
	
	\maketitle
	
	\textbf{Cognome e Nome:} \underline{\hspace{8cm}} \hfill
	\textbf{Data:} \underline{\hspace{3cm}}
	
	\vspace{0.5cm}
	
	\textbf{Durata:} 2 ore
	
	\vspace{0.5cm}
	
	\section*{Contesto}
	
	Un'azienda agricola desidera stimare la \textbf{resa agricola} di un terreno (in quintali per ettaro) in funzione di tre variabili:
	
	\begin{itemize}
		\item quantità di fertilizzante utilizzato (kg per ettaro),
		\item quantità di pioggia stagionale (mm),
		\item temperatura media stagionale (°C).
	\end{itemize}
	
	Si vuole costruire un modello di \textbf{regressione lineare multipla} utilizzando Python e la libreria \texttt{scikit-learn}.
	
Scrivere il codice in Python e rispondere alle domande di teoria:
	
	\begin{itemize}
		\item Esercizio 1
		\subitem generare dati simulati. 2 pt
		\subitem addestrare un modello di regressione, 2 pt
		\subitem utilizzare il modello per fare previsioni  1 pt.
		\subitem visualizzare una sezione del piano di regressione in un grafico 3D. 1 pt
		\item Esercizio 2
		\subitem Rispondere alle domande di teoria. 4 pt
	\end{itemize}
	
	\section*{Esercizio 1}
	
	Completa il seguente programma Python scrivendo il codice mancante a partire dai commenti.
	
	\begin{lstlisting}
		# Importare le librerie necessarie:
		# numpy
		# matplotlib.pyplot
		# LinearRegression da sklearn.linear_model
		# Axes3D da mpl_toolkits.mplot3d
		
		
		# -------------------------
		# 1) DATI (simulati)
		# -------------------------
		
		# fissare il seed del generatore casuale
		
		# definire il numero di osservazioni (n)
		
		# generare i valori casuali della variabile fertilizzante
		# (distribuzione uniforme tra 50 e 250)
		
		# generare i valori casuali della variabile pioggia
		# (distribuzione uniforme tra 200 e 700)
		
		# generare i valori casuali della variabile temperatura
		# (distribuzione uniforme tra 12 e 30)
		
		# generare il rumore casuale usando una distribuzione normale
		
		# definire la formula della resa agricola usando
		# fertilizzante, pioggia, temperatura e rumore
		
		
		# creare la matrice delle variabili indipendenti X
		
		# definire la variabile target y
		
		
		# -------------------------
		# 2) REGRESSIONE
		# -------------------------
		
		# creare il modello di regressione lineare
		
		# addestrare il modello con fit()
		
		# estrarre i coefficienti del modello
		
		# estrarre l'intercetta
		
		# calcolare il coefficiente di determinazione R^2
		
		# stampare l'equazione del modello e i parametri stimati
		
		
		# -------------------------
		# 3) PREVISIONI
		# -------------------------
		
		# definire due nuovi casi di input
		# (fertilizzante, pioggia, temperatura)
		
		# utilizzare il modello per effettuare le previsioni
		
		# stampare le previsioni in modo leggibile
		
		
		# -------------------------
		# 4) VISUALIZZAZIONE 3D
		# -------------------------
		
		# fissare un valore di temperatura
		
		# creare la figura e gli assi 3D
		
		# selezionare i punti con temperatura vicina al valore fissato
		
		# creare uno scatter plot dei punti selezionati
		
		# creare una griglia di valori per fertilizzante e pioggia
		
		# generare le matrici con meshgrid
		
		# calcolare i valori del piano di regressione
		
		# disegnare la superficie del piano
		
		# impostare le etichette degli assi
		
		# impostare il titolo del grafico
		
		# mostrare il grafico
	\end{lstlisting}
	
	\newpage
	
	\section*{Esercizio 2 - Domande teoriche}
	
	\begin{enumerate}
		
		\item Spiega il significato dell'istruzione che imposta il \textit{seed} del generatore casuale.
		
		\item Che cosa rappresentano i coefficienti del modello di regressione?
		
		\item Che cosa rappresenta l'intercetta del modello?
		
		\item Indica almeno due possibili variabili aggiuntive che potrebbero influenzare la resa agricola.
		
		\item Descrivi il funzionamento di un modello di regressione lineare
		
	\end{enumerate}
	
	

	
\end{document}